%==================================================================================================================
%==================================================== Différentes recombinaisons =======================================================
%==================================================================================================================

\begin{frame}[t]{Différentes types de recombinaisons}
  \vspace{-0.5em}
  \centering
      \begin{pccard}
        \centering
        \figrecombinaison
        \vspace{0.5em}
        {\scriptsize\bfseries\color{pcNavy}
        Différents types de recombinaisons : SRH, radiative et Auger}
      \end{pccard}
\end{frame}

%------------------------------------------------------------------------------------------------------------------

\begin{frame}[t]{Annexe : conditions aux limites — semi-conducteur}
  \vspace{-0.5em}

  \begin{columns}[T,onlytextwidth]

    \column{0.48\textwidth}
      \begin{pcbanner}{Contacts ohmiques (retenu)}
        \centering\scriptsize
        \textbf{Poisson : Dirichlet sur le potentiel}
        \[ \psi = \psi_0 + U_{\mathrm{app}} \]
        {\raggedright tension appliquée répartie symétriquement sur les deux contacts\par}
        \vspace{0.4em}
        \textbf{Densités : valeurs d'équilibre figées}
        \[ n = n_{\mathrm{eq}}, \qquad p = p_{\mathrm{eq}} \]
        {\raggedright nœuds de bord non mis à jour par le schéma temporel\par}
      \end{pcbanner}

    \column{0.48\textwidth}
      \begin{pcbanner}{Recombinaison de surface finie (perspective)}
        \begin{itemize}\footnotesize
          \item Hypothèse actuelle : contacts ohmiques idéaux (forte)
          \item Alternative : conditions de \textbf{Robin}
          \item Permet une recombinaison de surface finie aux contacts
          \item Objet de comparaison à venir
        \end{itemize}
      \end{pcbanner}

  \end{columns}

\end{frame}
%----------------------------------------------------------------------------------------

\begin{frame}[t]{Modèle 1D homogénéisé}
  \vspace{-0.5em}

  \begin{columns}[T,onlytextwidth]

    \column{0.46\textwidth}
      \begin{pccard}
        \centering
        \figschemaelectrochimique
        \vspace{0.4em}
        {\scriptsize\bfseries\color{pcNavy}
        Prise de moyenne volumique : 3D pore $\rightarrow$ 1D électrode}
      \end{pccard}

      \begin{pcbanner}{Principe}
        \begin{itemize}\footnotesize
          \item Moyenne sur un volume élémentaire représentatif
          \item Microstructure incorporée dans des coefficients effectifs
          \item Fermeture par vecteurs $b$, $d$, $e$
        \end{itemize}
      \end{pcbanner}

    \column{0.52\textwidth}
      \vspace{-0.5em}
      \begin{pcbanner}{Système homogénéisé}
        \centering\scriptsize
        \textbf{Conservation des espèces}
        \[ \varepsilon_l\frac{\partial \langle C_i\rangle^l}{\partial t}
           = \nabla\!\cdot\!\Big[ D_i^{\mathrm{tot}}\nabla\langle C_i\rangle^l
           + \tfrac{z_i F}{RT}\varepsilon_l D_i^{\mathrm{eff,mig}}
             \langle C_i\rangle^l \nabla\langle\varphi_l\rangle^l
           - \varepsilon_l \langle C_i\rangle^l \langle v\rangle^l \Big]
           + a_v^{\mathrm{cat}} R_i \]

        \textbf{Électroneutralité}
        \[ \sum_i z_i \langle C_i\rangle^l = 0 \]

        \textbf{Loi d'Ohm (phase solide)}
        \[ \nabla\!\cdot\!\big(\varepsilon_s \sigma^{\mathrm{eff}}
           \nabla\langle\varphi_s\rangle^s\big) = -\,a_v^{\mathrm{cat}}\, j \]

        \vspace{0.3em}
        \textbf{Coefficients effectifs}\\[0.2em]
        $D_i^{\mathrm{tot}} = \varepsilon_l D_i^{\mathrm{eff}} + D^{\mathrm{disp}}$,
        \quad $\sigma^{\mathrm{eff}}$, \quad $a_v^{\mathrm{cat}} = \theta a_v$
      \end{pcbanner}

  \end{columns}

\end{frame}

%----------------------------------------------------------------------------------------

\begin{frame}[t]{Annexe : réduction au module de Thiele}
  \vspace{-0.5em}

  \begin{columns}[T,onlytextwidth]

    \column{0.48\textwidth}
      \begin{pcbanner}{Hypothèses successives}
        \begin{itemize}\footnotesize
          \item Régime stationnaire
          \item Convection négligée ($\langle v\rangle^l = 0$)
          \item Migration négligée
          \item Potentiel solide uniforme ($\nabla\langle\varphi_s\rangle^s = 0$)
          \item Cinétique du 1\textsuperscript{er} ordre (Butler--Volmer linéarisée)
        \end{itemize}
      \end{pcbanner}

    \column{0.48\textwidth}
      \vspace{-0.5em}
      \begin{pcbanner}{Équation de diffusion--réaction}
        \centering\scriptsize
        \[ \frac{\mathrm{d}^2 \langle C_i\rangle^{l,*}}{\mathrm{d}\zeta^2}
           = \phi^2\, \langle C_i\rangle^{l,*} \]

        \textbf{Module de Thiele}
        \[ \phi^2 = -\,\frac{\nu_i\, a_v^{\mathrm{cat}}\, k(\eta)\, l_e^2}
                          {n_e F\, \varepsilon_l D_i^{\mathrm{eff}}} \]

        \textbf{Facteur d'efficacité}
        \[ \eta_E = \frac{\tanh(\phi)}{\phi} \]

        \vspace{0.2em}
        {\raggedright $\phi \ll 1$ : diffusion rapide, concentration uniforme\\
        $\phi \gg 1$ : réaction rapide, chute près de l'interface\par}
      \end{pcbanner}

  \end{columns}

\end{frame}

%--------------------------------------------------------------------------------

\begin{frame}[t]{Annexe : conditions aux limites — électrochimie}
  \vspace{-0.5em}

  \begin{columns}[T,onlytextwidth]

    \column{0.5\textwidth}
      \begin{pcbanner}{Système homogénéisé 1D}
        \centering\scriptsize
        \textbf{Collecteur de courant}
        \[ \langle N_i\rangle^l\!\cdot\!n = 0, \quad
           n\!\cdot\!\nabla\langle\varphi_l\rangle^l = 0, \quad
           \langle\varphi_s\rangle^s = E \]

        \vspace{0.3em}
        \textbf{Interface pore / milieu homogène}
        \[ \langle C_i\rangle^l = C_i^{\mathrm{hom}}, \quad
           \langle\varphi_l\rangle^l = \varphi_l^{\mathrm{hom}} \]
        \[ n\!\cdot\!\big(\varepsilon_s\sigma^{\mathrm{eff}}
           \nabla\langle\varphi_s\rangle^s\big) = 0 \]
      \end{pcbanner}

    \column{0.5\textwidth}
      \vspace{-0.5em}
      \begin{pcbanner}{Cas simplifié \ch{CO2 -> CO}}
        \centering\scriptsize
        \textbf{Concentration imposée au bulk ($\zeta = 0$)}
        \[ \langle C_{\ch{CO2}}\rangle^l(0) = C^{\mathrm{hom}}_{\ch{CO2}}, \quad
           \langle C_{\ch{CO}}\rangle^l(0) = C^{\mathrm{hom}}_{\ch{CO}} \]

        \vspace{0.3em}
        \textbf{Flux nul à la paroi ($\zeta = -1$)}
        \[ \frac{\mathrm{d}\langle C_{\ch{CO2}}\rangle^l}{\mathrm{d}\zeta}(-1) = 0, \quad
           \frac{\mathrm{d}\langle C_{\ch{CO}}\rangle^l}{\mathrm{d}\zeta}(-1) = 0 \]
      \end{pcbanner}

  \end{columns}

\end{frame}